\documentclass{article}
\usepackage[utf8]{inputenc}
\usepackage[T1]{fontenc}
\usepackage{geometry}
\usepackage{graphicx}
\usepackage{float}
\usepackage{hyperref}
\usepackage[english,indonesian]{babel}

\title{Optimasi Pemrosesan Template LaTeX Menggunakan Pendekatan Berbasis AI}
\author{Academic Support Agent}
\date{\today}

\begin{document}
\maketitle

\begin{abstract}

\end{abstract}

\chapter{PENDAHULUAN}


\chapter{TINJAUAN PUSTAKA}


% Original content (commented out):
% \documentclass[english]{article}
\usepackage{graphicx}
\usepackage{subfigure}
\usepackage[T1]{fontenc}
\usepackage[latin9]{inputenc}
\usepackage{babel}
\usepackage{amsmath}
\usepackage{amsthm}
\usepackage{amssymb}
\usepackage{amsfonts}
\usepackage{enumerate}
\usepackage{fancyhdr}
\usepackage{float}
\usepackage{url}
\usepackage{hyperref}
\usepackage{natbib}
\newcommand{\LaTeXs}{\LaTeX{\/} }
\newcommand{\matlabs}{{\sc Matlab{\/} }}
\newcommand{\matlab}{{\sc Matlab}}
\begin{document}

\title{Title}

\author{
  Author 1 \thanks{Put here comments or information about Author 1, if desired.} \and
  Author 2 \thanks{Put here comments or information about Author 2, if desired.} \and
  Author 3 \thanks{Put here comments or information about Author 3, if desired.}
}
% \date{xx/xx/xx} %uncomment this to place a specific date or leave it blank to erase the date 

\maketitle

\begin{abstract}
  Place your abstract here. 
\end{abstract}

\section{Introduction}\label{introduction}
The introduction goes here.  Start by drawing your reader in.  

Give the mathematical and scientific background and analysis required
to understand your project.  Include all relevant equations.

Tell us very briefly what you are doing in the project.
 
\section{Methods}\label{methods}
Methods goes here.  Methods are minimalistic.  Methods describe what
an intelligent reader capable of programming would need to know to
reproduce your work. You need to tell your reader you are using
Python, NumPy, SciPy, and Matplotlibs.  You need to mention that your
ODE-solver is \verb+odeint+.  However, the methods section is not a
Python programming guide!  Do not give a blow-by-blow of your code;
assume your reader can figure out how to do the line-by-line coding.

\section{USA population: results and conclusions}
Make sure your figures are readable when you upload your paper to
TurnItIn.  You can upload multiple PDFs to TurnItIn before the
deadline.  TurnItIn reduces the image quality of your document, often
making figures difficult to read.

\begin{figure}[H]
	\begin{centering}
  		\includegraphics[width=4in]{figure.png}
  		\caption{Place here information for the figure.\label{fig: floatfigure}
}
	\end{centering}
\end{figure}

You can reference the floating figure; it is figure \ref{fig: floatfigure}.  You
may find it weird that the label is in the caption, and it is.  This is a
\LaTeXs bug.  If it isn't in the caption, we will get the label for the section
and not for the figure.

Here's the table format for the Population Results:

\begin{tabular}{|c|c|c|c|c|c|c|}
\hline
Fit & \multicolumn{3}{|c|}{US population (millions)}&\multicolumn{3}{|c|}{Reduced Model Predictions}\\
\hline
Model & $P_0$ & $L$ & $r$ & 2010 Data & Prediction & Error \%\\
\hline
Exponential & $\qquad P_0   \qquad$ &  & $r$  & 308.7 & 341.4 & 10.6\\
\hline
Logistic & $P_0  $ & $\qquad L  \qquad$ & $r   $ & 308.7 &  & \\
\hline
Gompertz & $P_0 $& $L  $  & $r  $ & 308.7 & & \\
\hline
Fit & \multicolumn{3}{|c|}{Bryan TX population (thousands)}&\multicolumn{3}{|c|}{Reduced Model Predictions}\\
\hline
Exponential & $0.503$ &  & $0.0264$ & 65,660 & 76,394 & 16.4 \\
\hline
Logistic & $0.523$ & $10.17$ & $0.0421$ & 65,660 &  65,480 & -0.3\\
\hline
Gompertz & $0.512$ & $29.21$ & $0.0122$ & 65,660 & 68,036 & 3.6\\
\hline
\end{tabular}


\bigskip
And a table for comparing the different models:


\bigskip
\begin{tabular}{|c|c|c|c|c|c|c|}
\hline
&\multicolumn{3}{|c|}{$R^2$} & \multicolumn{3}{|c|}{Prediction Error \%}\\
\hline
Data & Exponential  & Logistic & Gompertz & Exponential & Logistic & Gompertz\\
\hline
USA &  &  & & 10.6\%  &  &  \\
\hline
Bryan, TX & 0.950 & 0.998 & 0.998  & 16.35\% & -0.27\% & 3.62\% \\
\hline
\end{tabular}

\section{Other population: results and conclusions}
Put other results and conclusions here.

\section{Conclusion}
Put the conclusions here. (\cite{taumanifesto})

\bibliographystyle{plainnat}
\bibliography{m475jmlinhart}%bibliography file is m475jmlinhart.bib!

\section*{Group Work Assessment}

Here's my group work assessment information.

This is what I did:
\begin{enumerate}\itemsep0pt
\item I did this.
\item And I did that.
\item I did another thing.
\end{enumerate}

\medskip
This is what my partner did:
\begin{enumerate}\itemsep0pt
\item They did this thing.
\item They did that thing.
\end{enumerate}

Tell me about any particular problems or successes you had? 

Is there anything you would do differently if you got a do-over?

Would you want to work with this person/these people again?

If you had  \$100 to distribute among the group members according to the contribution to the group project, how do you think this should be distributed?  Please give some reasons why.

\section*{Reflection}

Write your reflection as a letter, starting, ``Dear Dr. Linhart'' and
concluding Sincerely, (or Warmly, Cordially, Best, etc.) your name.
Your reflection should be at least 250 words long (this is 3/4 of a
page).  The main question is: what did you learn from this project?  I
really want to hear about the specifics and examples, not
generalities.  Here are some things you can put in your reflection:

\begin{enumerate}\itemsep0pt
\item What is the most important thing you got out of doing this
  project?  While the course is all about teaching you math,
  mathematical modeling, and writing/communicating, your answer still
  might be that you got a lesson on procrastination or working with
  others.  Know what it is and tell me about it.
\item This is a mathematical modeling class. Math and mathematical modeling are central
  learning objective for the course, so comment on these even if they
  are already familiar.  Some things to consider:
\begin{enumerate}[a)]\itemsep0pt
\item What math/mathematical models are in this project?
\item Anything new to you? Did you learn anything new about a familiar topic?
\item Did you see the math/mathematical model used in a way you haven't seen them used before?
\item Think back to the dialog you wrote at the beginning of the
  semester.  What steps in the mathematical modeling process do you
  see at work in this mathematical model?
\item Do these models work well or poorly, and why?  
\end{enumerate}
\item This is a writing and communication class. Since
  writing/communication is central to the course, comment on this.
  Some things to ponder:
\begin{enumerate}[a)]\itemsep0pt
\item  We are always emphasizing audience, purpose and logical organization.  Did these influence what/how you wrote?
\item Did you get any helpful feedback on a draft?
\item Have you done anything different in than
  what you've done before?  
\item Was writing the abstract easier this time?  
\end{enumerate}
\item What did you learn about Python in this project?
\item What did you learn about \LaTeXs in this project?
\item What did you learn about time management in this project?
\item What did you learn about working with others in this project?
\end{enumerate}
\end{document}


\end{document}